\documentclass[11pt,a4paper]{article}
\usepackage[margin=27mm]{geometry}
\usepackage{amsmath,amssymb,amsthm,hyperref}
\newtheorem{theorem}{Theorem}
\newtheorem{lemma}[theorem]{Lemma}
\newtheorem{corollary}[theorem]{Corollary}
\newcommand{\gen}[1]{\langle #1\rangle}
\newcommand{\F}{\mathbb F}
\newcommand{\Z}{\mathbb Z}
\title{Decorated points and tetrahedral geometries}
\author{Human Compression project}
\date{}
\begin{document}\maketitle

\begin{abstract}
The two large parabolics in the tetrahedral construction are stabilizers
of nonzero vectors modulo a subgroup of scalars. Their incidence is
therefore a determinant condition on decorated projective points.
This describes the relevant two types uniformly and exposes the
dependence on a choice of coordinates in each fibre. The same affine
coordinates give the full period lattices of the two toroidal residues.
We distinguish those lattices from the normal closure of a single
relation, which need not have the correct index. Two orbits of decorated
tetrahedra give all four types, their flags and their correlations.
The small characteristic is treated separately. The construction
recovers the two families of hypertopes while correcting their printed
geometric coordinates and some of their period labels.
\end{abstract}

\section{The parabolic groups}

Let $p\ne3$ be prime. Let $F=\F_p(e)$, where $e^2+e+1=0$ and
$e\ne1$, and put $q=|F|$. Thus $q=p$ when $p\equiv1\pmod3$,
and $q=p^2$ when $p\equiv2\pmod3$. Write
\[
 U(t)=\begin{pmatrix}1&t\\0&1\end{pmatrix},\qquad
 V(t)=\begin{pmatrix}1&0\\t&1\end{pmatrix},\qquad
 E=\begin{pmatrix}e&0\\0&e^2\end{pmatrix}.
\]
All matrix groups below are read in $G=\operatorname{PSL}_2(F)$;
the displayed matrices have determinant one. The source marking is
\[
 a_1=\begin{pmatrix}0&1\\-1&1\end{pmatrix},\quad
 a_2=\begin{pmatrix}e&e^2\\0&e^2\end{pmatrix},\quad
 a_3=\begin{pmatrix}e^2&e\\0&e\end{pmatrix}.
\]
For its parabolics $G_1=\gen{a_2,a_3}$ and
$G_0=\gen{a_1^{-1}a_2,a_1^{-1}a_3}$, direct multiplication gives
\[
 a_2=U(1)E,\quad a_3=U(1)E^{-1},\qquad
 a_1^{-1}a_2=V(1)E,\quad a_1^{-1}a_3=V(1)E^{-1}.
\]
In particular the lower-left entry of $a_1^{-1}a_3$ is $e^2$;
the isolated formula on source page~8 printing $e$ in that
position is a typographical error.

\begin{lemma}\label{groups}
The two marked parabolics are
\[
 G_1=U(F)\rtimes\gen E,\qquad G_0=V(F)\rtimes\gen E.
\]
They have order $3q$, and $G_0\cap G_1=\gen E$.
The three source generators generate $G$.
\end{lemma}
\begin{proof}
In either pair $X=T(1)E$, $Y=T(1)E^{-1}$, one has
$X^{-1}Y=E$ and $XE^{-1}=T(1)$. Conjugation by $E$ multiplies
the parameter of $U(t)$ by $e^2$ and that of $V(t)$ by $e$.
The additive spans of $1,e$ and of $1,e^2$ are both $F$.
Thus all translations occur, and no further elements occur since
the stated affine groups contain the generators. Their complements
have order three and meet their translation groups trivially.
An upper and a lower triangular matrix of the displayed forms can
agree projectively only when both off-diagonal entries vanish;
this proves the intersection statement.
Finally the generated group contains every $U(t)$ and $V(t)$.
Elementary row operations express every determinant-one matrix as
a product of these. If a leading pivot vanishes, the matrix
$U(1)V(-1)U(1)$ exchanges the two coordinate lines; otherwise
one eliminates the lower entry and then the upper entry.
Diagonal matrices are products as well, since
\[
 U(t)V(-t^{-1})U(t)
   =\begin{pmatrix}0&t\\-t^{-1}&0\end{pmatrix}\quad(t\ne0),
\]
and multiplying two such matrices gives any determinant-one diagonal
matrix. This proves generation without a subgroup classification.
\end{proof}

\section{Decorations and incidence}

Put $d=\gcd(2,q-1)$ and
$K=\{\pm1,\pm e,\pm e^2\}\le F^\times$, so $|K|=3d$.
A decorated point is an equivalence class
\[
             [v]_K=Kv,\qquad v\in F^2\setminus\{0\}.
\]
Forgetting the decoration gives its projective line $Fv$.
Each fibre has $(q-1)/(3d)$ elements and is a torsor for
$F^\times/K$: scalar multiplication acts freely and transitively.
There is no preferred element in a fibre.

\begin{theorem}\label{incidence}
The source cosets of types $0$ and $1$ are two copies of the
decorated points. Under the identifications
\[
 G_0g\longmapsto[g^{-1}(0,1)^T]_K,
 \qquad G_1g\longmapsto[g^{-1}(1,0)^T]_K,
\]
their incidence is
\[
                         \det(v,w)\in K.                 \tag{1}
\]
It is independent of the vector representatives. Each object has
$q$ neighbours, and this bipartite incidence graph is connected.
\end{theorem}
\begin{proof}
The stabilizer in $\operatorname{SL}_2(F)$ of $[(1,0)^T]_K$
consists of the upper triangular matrices
$\left(\begin{smallmatrix}k&b\\0&k^{-1}\end{smallmatrix}\right)$
with $k\in K$ and $b\in F$. Its projective image is exactly $G_1$.
The analogous lower triangular stabilizer is $G_0$.
The action on decorated points is transitive: a nonzero vector can
be completed to a matrix of determinant one. The orbit--coset
identifications follow, with the inverse accounting for the source's
right cosets.

If two cosets intersect, a common representative takes their two
base vectors to representatives of $[v]_K,[w]_K$. Their determinant
is then $1$ or $-1$ before multiplying by allowed scalar factors,
so (1) holds. Conversely, if $\det(v,w)\in K$, one can multiply
one vector by an element of $K$ to make the determinant one.
The resulting ordered pair is a matrix in $\operatorname{SL}_2(F)$
and supplies a common coset representative. Multiplying either
vector by a member of $K$ cannot change condition (1).

Fix $v=(1,0)^T$. Every adjacent class has a unique representative
of the form $(a,1)^T$, with $a\in F$. This proves the degree.
To prove connectedness, first take representatives $v,v'$ on
distinct projective lines in the same copy. The two independent
linear equations
\[
                  \det(v,w)=1,\qquad\det(v',w)=1
\]
have a solution $w\ne0$, which gives a path of length two between
their decorated classes. If two classes lie on the same projective
line, choose a third class on a different line and apply this argument
twice. All vertices in one copy are connected. Every vertex in the
other copy has a neighbour, so the whole graph is connected.
\end{proof}

Choose a vector $b_z$ on each projective line $z$. Writing its
decorated points as $[t b_z]_K$, with $tK\in F^\times/K$, gives
the explicit coordinate formula
\[
 [t b_z]_K * [u b_w]_K
 \quad\Longleftrightarrow\quad
             t u\det(b_z,b_w)\in K.                      \tag{2}
\]
Changing the chosen vector $b_z$ changes the scalar coordinate at
$z$ and the determinant factors together. This is the dependence
on the choice of representatives that a description by a product
set must retain.

Section~6, equation~(4) of the fixed source instead makes two
objects $(t,z)$ and $(u,w)$ incident exactly when $t=u$ and
$z\ne w$. For $m=(q-1)/(3d)>1$, this gives $m$ disjoint connected
components, each a complete bipartite graph on two copies of the
projective line with the matching deleted. The actual coset
truncation is connected by Theorem~\ref{incidence}; hence those
two incidence rules cannot be isomorphic. For example $q=13$
gives two components under the printed rule and one under (1).
The correction concerns the geometric description in Section~6;
it does not refute the separate coset construction of Section~4.

\section{Two orbits of decorated tetrahedra}

Write $x=[(1,0)^T]_K$, $y=[(0,1)^T]_K$ and $z=[(1,1)^T]_K$.
Put
\[
 T=\{x,y,z\},\qquad d_2=[(-e^2,1)^T]_K,\quad
 d_3=[(-e,1)^T]_K,\qquad D_i=T\cup\{d_i\}.
\]
The four points of either $D_i$ have different projective lines.
Every determinant of two distinct representatives in $D_i$ belongs
to $K$. For example the determinants involving $(1,1)^T$ and
$(-e^2,1)^T$ give $1+e^2=-e$; the other five determinants are
$\pm1$ or $\pm e^2$. The calculation for $D_3$ exchanges $e,e^2$.
Let $\mathcal D_i$ be the $G$-orbit of $D_i$.

\begin{lemma}\label{tetra-stabilizers}
The stabilizers of $D_2,D_3$ are respectively
$G_2=\gen{a_1,a_3}$ and $G_3=\gen{a_1,a_2}$, each acting as
$A_4$ on its four decorated points. Suppose $p$ is odd.
Then the stabilizer of $T$ is $\gen{a_1}\cong C_3$,
the two tetrahedron orbits are disjoint, and each triple in a
tetrahedron of one orbit belongs to exactly one tetrahedron of
the other orbit.
\end{lemma}
\begin{proof}
On either tetrahedron $a_1$ cycles $x,y,z$ and fixes its fourth
point. On $D_2$, $a_3$ fixes $x$ and cycles $y,z,d_2$;
on $D_3$, $a_2$ fixes $x$ and cycles $y,z,d_3$.
These assertions follow by multiplying the four column vectors;
all scalar factors are in $K$. The two displayed three-cycles
generate $A_4$. A projective linear transformation fixing three
different projective lines is the identity, so these actions are
faithful and both subgroups have order twelve.

Consider an element stabilizing $T$. If it fixes each of $x,y,z$,
a lift to $\operatorname{SL}_2(F)$ is diagonal, and fixing $z$
forces its diagonal entries to agree. It is scalar and hence the
identity in $G$. Thus the stabilizer embeds in $S_3$ and contains
the three-cycle $a_1$. An element transposing $x,y$ and fixing $z$
would have a lift
\[
                         \begin{pmatrix}0&k\\k&0\end{pmatrix},
                         \qquad k\in K,
\]
with determinant $-k^2=1$. For odd $p$, the square subgroup of
$K$ has order three and cannot contain $-1$. There is no such
transposition; composing any odd permutation with a power of
$a_1$ reduces to this case. The stabilizer is therefore $C_3$.

For $p$ odd, an element stabilizing $D_i$ can first be composed
with an element of the displayed $A_4$ to fix $d_i$.
It then stabilizes $T$ and already belongs to that $A_4$.
This proves the full stabilizer assertion. If a translate $gD_i$
contains $T$, transitivity of $G_i$ on the four triples of $D_i$
allows us to choose $h\in G_i$ with $ghT=T$. Consequently
$gD_i=ghD_i=D_i$. Thus $D_i$ is the unique member of its orbit
containing $T$. In particular $D_2,D_3$ are in different orbits.
Transporting this argument by $G$ proves the assertion for every
triple in either orbit.

The remaining characteristic is $p=2$, so $q=4$ and $K=F^\times$.
The decorated points are the five ordinary projective points.
The group $G$ has order sixty and acts faithfully on them, with
the three generators acting by even permutations. It is $A_5$.
Both $D_i$ are four-subsets, and their orbits are the five
four-subsets. Their stabilizers have order twelve and are precisely
the already exhibited $G_i$. The assertions specific to odd $p$
are not used in this case.
\end{proof}

Define a four-type incidence system as follows. Types $0,1$ are
two copies of the decorated points; types $2,3$ are the orbits
$\mathcal D_2,\mathcal D_3$, retaining the type even if the orbits
coincide. Between different types use
\[
\begin{array}{c|l}
\text{types}&\text{incidence}\\\hline
0,1&\det(v,w)\in K,\\
i,j\ (i<2\le j)&[v]_K\in D,\\
2,3&|D\cap D'|=3.
\end{array}                                                     \tag{3}
\]

\begin{theorem}\label{full-geometry}
The incidence system (3) is the source's right-coset system
$\Gamma(G;(G_i)_{i=0}^3)$. It is a hypertope with tetrahedral
diagram having all six bonds labelled three. The group $G$ has
two free orbits on its chambers; adjacent chambers belong to
different orbits. All rank-three flags form one free $G$-orbit
for each choice of three types.
\end{theorem}
\begin{proof}
The point stabilizers and the first incidence rule were proved in
Theorem~\ref{incidence}. Lemma~\ref{tetra-stabilizers} identifies
the other two stabilizers. Thus the object map sends a right coset
$G_i g$ to $g^{-1}$ applied to the corresponding base object
$y,x,D_2,D_3$. Every incident point--tetrahedron pair is in the
orbit of the corresponding base pair, because $A_4$ is transitive
on the four vertices. For odd $p$, every tetrahedron of type $2$
has exactly four neighbours of type $3$, one for each triple.
Its $A_4$ stabilizer is transitive on these triples. Hence these
pairs too form exactly the orbit of the base pair $(D_2,D_3)$.
Pair incidence in a coset system is exactly the orbit of its base
pair, proving the claim about all six relations.

For the moment suppose $p$ is odd. A flag of types $012$ or $013$
is a tetrahedron with two ordered distinct vertices. Its stabilizer
$A_4$ acts regularly on the twelve ordered pairs of distinct
vertices. These rank-three flags therefore form one free orbit.
A flag of types $023$ or $123$ consists of adjacent tetrahedra
and a point of their common triple. The stabilizer of the pair is
$C_3$, acting regularly on that triple. These flags likewise form
one free orbit.

There are exactly two extensions of each such flag to a chamber.
For types $023$ or $123$, the missing point can be either of the
other two vertices in the common triple. For types $012$ or
$013$, the missing tetrahedron is specified by a triple of the
given tetrahedron containing the two given points. There are
exactly two choices, giving different opposite tetrahedra.
Every incident pair extends to one of these flags: for a point
and a tetrahedron choose another vertex; for adjacent tetrahedra
choose a point in their common triple; and the point pair case
follows by transporting the base point pair. Every object has
an incident pair. Thus this is a geometry, and the two-extension
calculation proves thinness.

The stabilizer of each base incident pair is cyclic of order
three. For pair types $01$, $02$, $03$, $12$, $13$ and $23$,
these groups are respectively
\[
 \gen E,\quad\gen{a_1^{-1}a_3},\quad\gen{a_1^{-1}a_2},\quad
 \gen{a_3},\quad\gen{a_2},\quad\gen{a_1}.
\]
For point--tetrahedron pairs this follows from the $A_4$ action;
the other two cases were proved above. Since triple stabilizers
are trivial, every rank-two residue has three objects of each
of its two types. Thinness makes its graph a simple bipartite
graph of degree two, hence a six-cycle. This proves both
connectedness of these residues and the bond labels.

Each rank-three residue is the coset geometry inside its object
stabilizer $G_i$ with parabolics $G_i\cap G_j$. This assertion
includes incidence: a pair in such a residue, together with the
fixed object, is a rank-three flag, so the transitivity just
proved puts it in the orbit of the base flag under $G_i$.
These three cyclic parabolics generate $G_i$ in every case, as
the displayed generators show. Their coset incidence graph is
connected: writing a group element as a product of elements of
these parabolics gives a path between the corresponding chambers
and their objects. The whole geometry is connected already on
types $0,1$, and every tetrahedron has point neighbours. All
required residues are therefore connected.

There are exactly $2|G|$ chambers, because each rank-three flag
has two extensions and there are $|G|$ flags of any fixed three
types. Chamber stabilizers are trivial. If adjacent chambers
were in the same orbit, a group element taking one to the other
would fix their common rank-three flag and would be the identity,
a contradiction. This proves the orbit assertions.

For $q=4$, identify each tetrahedron with its missing projective
point. All four types are then copies of the same five-element
set, with incidence meaning unequal labels. A chamber is an
ordered choice of four different labels. Every three-label flag
has two extensions; every rank-two residue is a six-cycle, and
every rank-three residue is three copies of four labels with
unequal-label incidence, which is connected. The $A_5$ action
on ordered triples of different labels is free and transitive.
These observations give all the same conclusions directly,
including the coset identification by the stabilizers.
\end{proof}

The six pair stabilizers are distinct, and their intersections
are trivial whenever a third base object is fixed. Together with
$G_i\cap G_j$ as displayed, this proves the full intersection
property for the rotation marking: a subgroup indexed by two
labels is one of these cyclic groups, one indexed by three is a
maximal parabolic, and the intersections are their indexed meets.
For example intersecting a cyclic group with a maximal parabolic
that does not contain it fixes three base objects and is trivial.
Two distinct cyclic groups have trivial intersection by the same
argument. The generators are independent since omitting one
leaves a proper maximal parabolic. This also recovers the source's
$C^+$-group claims, rather than assuming them to construct the
geometry.

\section{Residues and symmetries}

The residue of a tetrahedron is three copies of its four vertices:
the opposite tetrahedron is labelled by the one missing vertex
of its common triple. Incidence again means unequal labels.
Its rotation group is $A_4$; the two triangular translation words
are distinct nonidentity double transpositions and generate its
Klein four subgroup. Their squares are the identity, so the
translation kernel contains $2\Z^2$. Its index is four because
the image is that Klein four group. The kernel is therefore
exactly $2\Z^2$, giving the marked hypermap $(3,3,3)_{(0,2)}$.

The residue of a point has rotation group $G_0$ or $G_1$, with
the ordered generators used in the next section. Its actual
incidence is the internal coset incidence by the rank-three flag
argument in Theorem~\ref{full-geometry}. The classical triangular
rotation group maps onto it with a translation kernel. Therefore
the residue is the triangular toroidal quotient by precisely that
kernel, not merely by a relation satisfied in the finite group.
In either point residue the three cyclic parabolics, in order,
are $\gen{X^{-1}Y}$, $\gen Y$ and $\gen X$; they have order
three and trivial pairwise intersections. Thus this is the
standard marked $C^+$ coset construction. The residue stabilizer
has two free chamber orbits, inherited from the two global orbits,
and adjacent chambers are in opposite orbits.
The full kernels and the resulting correction of the printed
period parameters are calculated below.

Consider the two linear maps
\[
       A=\begin{pmatrix}-1&1\\0&1\end{pmatrix},\qquad
       B=\begin{pmatrix}0&1\\1&0\end{pmatrix}.
\]
Both have determinant $-1\in K$ and normalize $\operatorname{SL}_2(F)$.
They act on decorated points and tetrahedra, preserve all the
relations (3), and exchange $D_2,D_3$. The latter assertion
follows by applying each map to their four displayed vectors.
Thus $A$, accompanied by the type permutation $(2\ 3)$, is a
correlation $\mu$. It sends the base chamber
$(y,x,D_2,D_3)$ to $(z,x,D_2,D_3)$, its type-zero neighbour.
The map $B$, accompanied by $(0\ 1)(2\ 3)$, is a correlation
$\eta$ fixing the base chamber as a set. Both normalize the
$G$ action. Accordingly $\mu$ exchanges its two chamber orbits,
whereas $\eta$ preserves them. This proves the improper and
proper correlation assertions by the same object model.

When $q=p^2$, Frobenius $\sigma:t\mapsto t^p$ exchanges $e,e^2$.
On the object model it preserves determinant membership in $K$,
fixes $x,y,z$ and exchanges $D_2,D_3$. Accompanied by the type
permutation $(2\ 3)$ it is a correlation fixing the base chamber.
Its composition with $\mu$ is type preserving and takes the base
chamber to its type-zero neighbour. Together with $G$ it acts
transitively on all chambers. The hypertope is regular.

For $q=p$, the kernel calculation below shows that neither
marked point residue admits an automorphism inverting its two
standard rotation generators. The acquired theorem on two chamber
orbits~\cite{reflection}, applied with the preceding actual
coset and chamber identifications, says that the residue is
regular exactly when this inversion descends. Its proof labels
all chambers $(g,+),(g,-)$ and uses the explicit reflection
$(g,\pm)\mapsto(\theta(g),\mp)$ for an inverting automorphism
$\theta$; the converse follows by conjugating the rotations
by a normalized chamber reflection. The non-invariant kernel
therefore makes the point residue chiral.
A regular geometry has regular residues:
chamber transitivity restricts to chamber transitivity in the
stabilizer of the fixed flag. Hence the full geometry is not
regular. Its two $G$-orbits and thin residual connectedness show
that it is chiral, with full type-preserving group $G$.
Here we use the elementary gallery fact that in a thin residually
connected geometry any type-preserving automorphism fixing a
chamber is the identity: propagate across its uniquely determined
adjacent chambers. Residual connectedness gives gallery
connectedness by induction on rank. In particular the full group
cannot split either transitive $G$-orbit or contain an additional
element preserving both orbits.

For $q=4$ the full correlation group is $S_5\times S_4$.
Indeed the common permutation of the five labels and any
permutation of the four types preserve unequal-label incidence.
Conversely, remove the induced type permutation from a correlation.
For two different types, the unique non-neighbour of a label is
its equal label in the other copy. The permutations on the four
copies must therefore agree. The type-preserving group is $S_5$,
and the given rotation group is $A_5=\operatorname{PSL}_2(4)$.
This corrects the final isomorphism in the source's first remark
on page~19: $S_5\times S_4$ is not
$\operatorname{PSL}_2(4)\times S_4$.


\section{The complete translation kernels}

Use the source's ordered generators $X,Y$ for either $G_1$ or
$G_0$. In the standard triangular translation coordinates put
\[
                        u=XY^{-1}X,\qquad v=XY.
\]
The triangular rotation presentation and its translation lattice
$\Z^2$ are retained classical inputs, already developed in the
acquired hexagonal construction. The preceding matrices give
\[
\begin{array}{c|cc|c}
 &u&v&\text{image of }u^a v^b\\\hline
 G_1&U(1)&U(-e)&U(a-eb)\\
 G_0&V(1)&V(-e^2)&V(a-e^2b).
\end{array}                                                   \tag{4}
\]
The complement still has order three, so the kernel of either
marked triangular presentation is wholly in its translation subgroup.

\begin{theorem}\label{periods}
If $q=p^2$, the period lattice of either residue group is
$p\Z^2$. If $q=p$, choose integer representatives $e_0,e_1$ of
$e,e^2\in\F_p$. The two period lattices are respectively
\[
 L_1=\gen{(p,0),(e_0,1)}_{\Z},\qquad
 L_0=\gen{(p,0),(e_1,1)}_{\Z}.                              \tag{5}
\]
Both have index $p$. In contrast, the normal translation lattice
generated by the single source vector $(e_0,1)$ has index
$e_0^2+e_0+1$, which can be larger than $p$.
\end{theorem}
\begin{proof}
When $e\notin\F_p$, the elements $1,e$ are linearly independent
over $\F_p$, as are $1,e^2$. Equation (4) therefore has kernel
exactly $p\Z^2$. When $e\in\F_p$, the equation
$a-e_0b\equiv0\pmod p$ is equivalent to
$(a,b)=c(p,0)+b(e_0,1)$. This proves the first lattice in (5);
the other is identical with $e_1$ substituted. Their determinants
are $p$.

Conjugation by the triangular order-three complement acts on
translation coordinates by
\[
                R(a,b)=(-a-b,a),\qquad R^2+R+I=0.
\]
One may replace $R$ by $R^{-1}$ without changing an orbit lattice.
The normal closure of a translation vector $(a,b)$ is its orbit
span. The two columns $(a,b)$ and $R(a,b)$ have determinant
$a^2+ab+b^2$. Thus the orbit span of $(e_0,1)$ has index
$e_0^2+e_0+1$. It is contained in the actual kernel by (4),
but equality requires equality of these indices.
\end{proof}

For a concrete discrepancy take $p=19$ and $e=7$; then
$7^2+7+1=57$. The actual parabolic group has order $3p=57$,
whereas the single-period presentation printed as $(3,3,3)_{(7,1)}$
has order $3\cdot57=171$. The other printed vector $(1,7)$ has
the same norm. The missing condition is the second lattice
generator in (5). In the $G_1$ coordinates a suitable single
orbit generator is $(2,3)$: its norm is $19$, and
$2-7\cdot3\equiv0\pmod{19}$. Its orbit span therefore equals
$L_1$. In the $G_0$ coordinates $(3,2)$ works, since $e^2=11$
and $3-11\cdot2\equiv0\pmod{19}$.

The failure of the printed period vector does not invalidate the
chirality argument. To make this precise without choosing a norm
representation, an automorphism inverting both $X,Y$ would send
\[
                        u\longmapsto u^{-1},\qquad
                        v\longmapsto u^{-1}v.
\]
These identities follow directly from the matrices, or from the
triangular affine presentation. It would therefore preserve the
kernel under $(a,b)\mapsto(-a-b,b)$. For $q=p$, this map takes
$(e_0,1)\in L_1$ to $(e_1,1)\notin L_1$, since $e^2\ne e$.
It similarly fails to preserve $L_0$. Thus no such automorphism
exists for either marked residue group. When $q=p^2$, the lattice
$p\Z^2$ is preserved. These are the exact kernel statements needed
for the residue symmetry test. Section~4 identified the actual
residue geometry with the corresponding toroidal quotient.

\section{Interfaces with the source's intermediate lemmas}

The general coset construction and its transitivity on incident
pairs are the acquired Tits construction. Its right action need
not be faithful for arbitrary input subgroups; faithfulness here
follows from the trivial chamber stabilizer. Source Proposition~2
also has a short general proof. With at least two rotation
generators, every generator belongs to one of the maximal
parabolics, so their union generates the group. Chambers obtained
by multiplying successive such factors share an object. Every
object belongs to a chamber, giving a connected incidence graph.

The coset transversals in source Lemma~9 can also be read from
the object model. For $j=0,1$, the translation group $V(F)$ or
$U(F)$ is a transversal of each cyclic subgroup
$G_i\cap G_j$ in $G_j$: it has order $q$, the intersection
with a subgroup of order three is trivial, and $|G_j|=3q$.
For $j=2,3$, the displayed three-cycles on $D_j$ show that the
source's four listed representatives send the relevant fixed
vertex to its four different vertices. They are therefore
transversals too. This proves the literal neighbour lists of
Lemma~9. Source Proposition~10 is now an immediate consequence
of transitivity on rank-three flags: an object incident to two
base objects lies in their common stabilizer orbit.

For completeness, the identity behind source Lemma~18 is
\[
 A V(c) A a_1
   =\begin{pmatrix}-c&1\\-c-1&1\end{pmatrix}
   =U(1)V(-c-1)=a_1V(-c).
\]
Because $U(1)\in G_1$ and $a_1\in G_2\cap G_3$, and both
$c\mapsto-c-1$ and $c\mapsto-c$ permute $F$, it gives
exactly the three stated coset-set equalities. This calculation
is no longer needed to prove that $\mu$ preserves incidence:
the determinant and membership formulas prove that at once.

\section{Attribution and scope}

The field family and three matrices are those of Catalano,
Fernandes, Hubard and Leemans~\cite{source}. The affine descriptions
of the parabolics are their Lemma~6, here derived from two common
identities. Vector stabilizers, determinant coordinates and the
triangular translation group are standard ingredients. The addition
is the common model of decorated points and tetrahedra, which
derives the flag and correlation assertions from their small
stabilizers. The incidence formula and complete kernel calculation
also identify two discrepancies in the fixed edition. The chamber
inversion theorem is acquired from the existing common theory,
not reproved for this family. No historical priority is claimed
for these methods.
The four-type construction and its proper and improper correlations
follow from the same decorated tetrahedra. The complete source
disposition and independent review are recorded separately; this
working manuscript does not itself certify their completion.
It gives no new conjecture answer or compression measurement.

The second remark on source page~19 states that the finite groups
$G^+=\operatorname{PSL}_2(q)$ are index-four subgroups of the
infinite hyperbolic Coxeter group $[6,3,6]$. Read literally this
is impossible: a group with a finite subgroup of finite index
is finite. A construction using a universal subgroup and then
finite quotients would require a different statement. We do not
silently substitute that possible interpretation for the printed one.

\begin{thebibliography}{9}
\bibitem{source} D.~Catalano, M.~E.~Fernandes, I.~Hubard and
D.~Leemans, \emph{Hypertopes with tetrahedral diagram},
Electronic Journal of Combinatorics 25(3) (2018), P3.22.
Fixed published edition, catalogue L041, Lemmas~6 and~14,
Section~6, equation~(4).
\bibitem{acquired} Human Compression project,
\emph{Hexagonal extensions of toroidal maps and hypermaps},
acquired triangular rotation and translation presentations.
\bibitem{reflection} Human Compression project,
\emph{Residue stabilizers and distinguished generators}, Section~6,
and \emph{Reconstructing a reflection from two chamber orbits},
accepted common theory, topic-chamber/v2, Section~2.
The latter makes explicit the criterion of Fernandes, Leemans
and Ivi\'c Weiss, \emph{Highly symmetric hypertopes},
Aequationes Mathematicae 90 (2016), Theorem~8.2.
\end{thebibliography}
\end{document}
